% Part: proof-theory
% Chapter: normalization
% Section: reductions

\documentclass[../../../include/open-logic-section]{subfiles}

\begin{document}

\olfileid{pt}{nor}{red}

\olsection{تحويلات الاختزال}

إذا كان طول قطع في~!!a{proof} ضمن~\Log{N1i} يساوي~$1$، فهو مقطع يتألف من
ظهور واحد لـ~!!{formula}، وتكون تلك~!!{formula} نتيجةً إما لاستدلال
!!a{introduction} أو لـ~\FalseInt، ومقدمةً كبرى لاستدلال~!!a{elimination}
في الوقت نفسه. وفي برهان التطبيع ``تُختزل'' هذه القطوع باستبدال~!!{proof}
الجزئي الذي ينتهي بمثل هذا الزوج من~!!{introduction}/!!{elimination}
بـ~!!{proof} جديد حُذف منه زوج الاستدلالين. وقد ينتج عن هذا الاستبدال قطوع
جديدة، غير أنه يمكن التحقق من أن رتبتها أدنى من رتبة القطع المختزل. ولذلك،
عند تطبيق التحويل على قطع أقصى مختار علويًا وفي أقصى اليمين، يكون~!!{proof}
الجديد إما أقصر في طول القطع (لأن قطعًا أقصى طوله~$1$ قد حُذف) وإما تكون
رتبة القطع في~!!{proof} قد انخفضت (إذا كان ذلك هو القطع الأقصى الوحيد).

فمثلًا، إذا كان الاستدلالان المعنيان هما~\Intro{\lif} و\Elim{\lif}، وكانت
$!A \ident !B \lif !C$، وكان هناك~!!{proof} جزئي~$\delta_1$ ينتهي هكذا:
\[
\AxiomC{$\Discharge{!B}{x}$}
\RightLabel{$\delta_2$}
\DeduceC{$!C$}
\DischargeRule{\Intro{\lif}}{x}
\UnaryInfC{$!B \lif !C$}
\AxiomC{}
\RightLabel{$\delta_3$}
\DeduceC{$!B$}
\RightLabel{\Elim{\lif}}
\BinaryInfC{$!C$}
\DisplayProof
\]
فإننا نحذف هذا القطع باستبدال البرهان الجزئي~$\delta_1$ داخل~$\delta$
بما يأتي:
\[
\AxiomC{}
\RightLabel{$\delta_3$}
\DeduceC{$!B$}
\RightLabel{$\Subst{\delta_2}{\delta_3}{!B^x}$}
\DeduceC{$!C$}
\DisplayProof
\]
ونعني بـ~$\Subst{\delta_2}{\delta_3}{!B^x}$ الـ~!!{proof} الناتج من
$\delta_2$ باستبدال نسخة كاملة من~!!{proof}~$\delta_3$ بكل فرض~$!B$ موسوم
بـ~$x$. ونعيد تسمية المتغيرات الذاتية والمقيّدة في كل نسخة من~!!{proof}
$\delta_3$ بصورة حديثة ومتباينة زوجيًا تجنبًا للأسر، مع إبقاء وسوم الفروض
المفتوحة وصيغها نفسها، وإن جاز أن يزداد عدد وقائعها. ولا يبقى في البرهان
الجديد فرض مفتوح من الصورة~$!B$ موسوم بـ~$x$. كما لا توجد فروض أخرى موسومة
بـ~$x$، لأننا نفترض كالمعتاد أن~!!{formula}s مختلفة لا تقع فروضًا بالوسم
نفسه. ومن ثم يكون لـ~!!{proof} الجديد الفروض المفتوحة نفسها التي لـ
$\delta_1$؛ وكل استدلال في~$\delta_2$ يفرّغ فروضًا يفرّغ الفروض نفسها في
$\Subst{\delta_2}{\delta_3}{!B^x}$، وكل استدلال في~$\delta_3$ يفرّغ فروضًا
قد تظهر منه نسخ متعددة تفرّغ الوقائع المقابلة.

كل قطع يقع كله داخل~$\delta_2$ أو كله خارج~$\delta_1$ في~$\delta$ لا
يتغير: فهو يشتمل على~!!{formula}s نفسها ويمر بالاستدلالات نفسها، فتظل رتبته
وطوله كما كانا. أما القطع الواقع داخل~$\delta_3$ فقد يُنسخ؛ ولكل نسخة الرتبة
والطول نفسيهما، وتتحكم استراتيجية اختيار القطع الأقصى العلوي وفي أقصى
اليمين في تعدديته. وقد حذفنا قطعًا أقصى طوله~$1$، لكن يلزم بحث أمرين قبل
استنتاج انخفاض مقياس~$\delta$.

غير أن أمرين قد يحدثان:
\begin{enumerate}
  \item باستبدال~$\delta_3$ بكل فرض~$!B^x$، نكون قد ضاعفنا جميع القطوع في
  $\delta_3$. فإذا كان بعضها أقصى، فقد يزيد طول القطع في~!!{proof} الجديد
  على طول القطع في~!!{proof} القديم. ولمنع ذلك نطبّق تحويلات الاختزال فقط
  على قطوع قصوى \emph{علوية وفي أقصى اليمين} معًا؛ أي لا يوجد قطع أقصى فوق
  القطع المختزل في برهانه الجزئي، ولا على فرع في~$\delta$ إلى يمينه. فلو
  وُجد قطع أقصى في~$\delta_3$ لكان إلى يمين القطع المختزل هنا، ولو وُجد في
  الجزء العلوي الآخر لانتُهكت العلوية. ويضمن هذا الاختيار أننا إما نقلل طول
  القطع في~!!{proof} أو نقلل رتبته.
  \item إذا كان فرض~$!B^x$ في~$\delta_2$ المقدمة الكبرى لاستدلال
  !!a{elimination}، وكان آخر استدلال في~$\delta_3$ هو~\Elim{\lor} أو
  \Elim{\lexists} أو استدلال~!!a{introduction} أو~\FalseInt، فقد أنشأنا
  قطعًا جديدًا بوصل~$\delta_3$ بـ~$\delta_2$ عند ذلك الفرض. فما كان مجرد
  فرض في~$\delta_2$ صار قطعًا طوله~$1$ (نتيجة~!!a{introduction} أو
  \FalseInt ومقدمةً كبرى لاستدلال~!!a{elimination})، أو صار آخر ظهور
  لـ~!!{formula} في مقطع قطع (أي مقدمةً كبرى لاستدلال~!!a{elimination}
  ونتيجةً لـ~\Elim{\lor} أو~\Elim{\lexists}). وإذا كانت~$B^x$ مقدمةً
  صغرى لـ~\Elim{\lor} أو~\Elim{\lexists}، فقد كانت أصلًا أول ظهور
  لـ~!!{formula} في مقطع، وربما في مقطع قطع. وقد يطول ذلك المقطع في
  !!{proof} الجديد. لكن لاحظ أن~!!{formula} التي تكوّن هذا القطع الجديد أو
  مقطع القطع القائم هي~$!B$، وأن
  $\depth{!B} < \depth{!B \lif !C}$. لذلك، مع أننا ربما أضفنا قطوعًا أو
  زدنا أطوالها، فلا شيء منها قطع أقصى. ومن ثم إما أننا خفضنا رتبة القطع، أو
  خفضنا طول القطع إن بقيت قطوع بالرتبة نفسها.
\end{enumerate}

تسمى عملية اختزال قطع طوله~$1$ \emph{تحويل اختزال} (أو \emph{تحويل
التفاف}). وترد أنماط بعض أزواج القواعد في
\olref{tab:reduction-conversions}. (وتُترك الأنماط الباقية تمارين.)

\begin{table}
\begin{multline*}
\bottomAlignProof
\AxiomC{$\Discharge{!B}{x}$}
\RightLabel{$\delta_2$}
\shortDeduce
\DeduceC{$!C$}
\DischargeRule{\Intro{\lif}}{x}
\UnaryInfC{$!B \lif !C$}
\AxiomC{}
\RightLabel{$\delta_3$}
\shortDeduce
\DeduceC{$!B$}
\RightLabel{\Elim{\lif}}
\BinaryInfC{$!C$}
\DisplayProof
\quad \longrightarrow\quad
\shoveright{\bottomAlignProof
\AxiomC{}
\RightLabel{$\delta_3$}
\shortDeduce
\DeduceC{$!B$}
\RightLabel{$\Subst{\delta_2}{\delta_3}{!B^x}$}
\shortDeduce
\DeduceC{$!C$}
\DisplayProof}
\\[4ex]
\shoveleft{\bottomAlignProof
\AxiomC{}
\RightLabel{$\delta_2$}
\shortDeduce
\DeduceC{$!B$}
\AxiomC{}
\RightLabel{$\delta_3$}
\shortDeduce
\DeduceC{$!C$}
\RightLabel{\Intro{\land}}
\BinaryInfC{$!B \land !C$}
\RightLabel{\Elim{\land}[1]}
\UnaryInfC{$!B$}
\DisplayProof
\quad \longrightarrow\quad
\bottomAlignProof
\AxiomC{}
\RightLabel{$\delta_2$}
\shortDeduce
\DeduceC{$!B$}
\DisplayProof}\qquad
\shoveright{\bottomAlignProof
\AxiomC{}
\RightLabel{$\delta_2$}
\shortDeduce
\DeduceC{$!B$}
\AxiomC{}
\RightLabel{$\delta_3$}
\shortDeduce
\DeduceC{$!C$}
\RightLabel{\Intro{\land}}
\BinaryInfC{$!B \land !C$}
\RightLabel{\Elim{\land}[2]}
\UnaryInfC{$!C$}
\DisplayProof
\quad \longrightarrow\quad
\bottomAlignProof
\AxiomC{}
\RightLabel{$\delta_3$}
\shortDeduce
\DeduceC{$!C$}
\DisplayProof}
\\[4ex]
\bottomAlignProof
\AxiomC{}
\RightLabel{$\delta_2$}
\shortDeduce
\DeduceC{$!B$}
\RightLabel{\Intro{\lor}[1]}
\UnaryInfC{$!B \lor !C$}
\AxiomC{$\Discharge{!B}{x}$}
\RightLabel{$\delta_3$}
\shortDeduce
\DeduceC{$!D$}
\AxiomC{$\Discharge{!C}{y}$}
\RightLabel{$\delta_4$}
\shortDeduce
\DeduceC{$!D$}
\DischargeRule{\Elim{\lor}}{x\,y}
\TrinaryInfC{$!D$}
\DisplayProof
\quad \longrightarrow\quad
\shoveright{\bottomAlignProof
\AxiomC{}
\RightLabel{$\delta_2$}
\shortDeduce
\DeduceC{$!B$}
\RightLabel{$\Subst{\delta_3}{\delta_2}{!B^x}$}
\shortDeduce
\DeduceC{$!D$}
\DisplayProof}
\\[2ex]
\shoveleft{\bottomAlignProof
\AxiomC{}
\RightLabel{$\delta_2$}
\shortDeduce
\DeduceC{$!B(c)$}
\RightLabel{\Intro{\lforall}}
\UnaryInfC{$\lforall[x][!B(x)]$}
\RightLabel{\Elim{\lforall}}
\UnaryInfC{$!B(t)$}
\DisplayProof
\quad \longrightarrow\quad
\bottomAlignProof
\AxiomC{}
\RightLabel{$\Subst{\delta_2}{t}{c}$}
\shortDeduce
\DeduceC{$!B(t)$}
\DisplayProof}
\\[4ex]
\shoveright{\bottomAlignProof
\AxiomC{}
\RightLabel{$\delta_2$}
\shortDeduce
\DeduceC{$!B(t)$}
\RightLabel{\Intro{\lexists}}
\UnaryInfC{$\lexists[x][!B(x)]$}
\AxiomC{$\Discharge{!B(c)}{x}$}
\RightLabel{$\delta_3$}
\shortDeduce
\DeduceC{$!D$}
\DischargeRule{\Elim{\lexists}}{x}
\BinaryInfC{$!D$}
\DisplayProof
\quad \longrightarrow\quad
\bottomAlignProof
\AxiomC{}
\RightLabel{$\delta_2$}
\shortDeduce
\DeduceC{$!B(t)$}
\RightLabel{$\Subst{\Subst{\delta_3}{t}{c}}{\delta_2}{!B(t)^x}$}
\shortDeduce
\DeduceC{$!D$}
\DisplayProof}
\end{multline*}
\caption{بعض تحويلات الاختزال في~\Log{N1i}.}
\ollabel{tab:reduction-conversions}
\end{table}

\begin{prob}
اكتب تحويل الاختزال لـ~\Intro{\lnot} متبوعةً بـ~\Elim{\lnot}.
\end{prob}

تبقى حالة واحدة للنظر. يشمل تعريف القطع الحالة التي يكون فيها طول القطع
$1$ وتكون~$!A$ نتيجة~\FalseInt{} ومقدمةً كبرى لاستدلال~!!a{elimination}
في الوقت نفسه. وتُعالج هذه الحالة بطريقة شبيهة بالحالة التي تكون فيها~$!A$
نتيجة قاعدة~!!a{introduction}. فمثلًا، استبدل
\[
\bottomAlignProof
\AxiomC{}
\RightLabel{$\delta_2$}
\DeduceC{$\lfalse$}
\RightLabel{\FalseInt}
\UnaryInfC{$!B \lif !C$}
\AxiomC{}
\RightLabel{$\delta_3$}
\DeduceC{$!B$}
\RightLabel{\Elim{\lif}}
\BinaryInfC{$!C$}
\DisplayProof
\quad\text{بـ}\quad
\bottomAlignProof
\AxiomC{}
\RightLabel{$\delta_2$}
\DeduceC{$\lfalse$}
\RightLabel{\FalseInt}
\UnaryInfC{$!C$}
\DisplayProof
\]
وينبغي بعض الحذر في حالتي~$!A \ident (!B \lor !C)$ و
$!A \ident \lexists[x][!B(x)]$. فمثلًا، لو استبدلنا
\[
\bottomAlignProof
\AxiomC{}
\RightLabel{$\delta_2$}
\DeduceC{$\lfalse$}
\RightLabel{\FalseInt}
\UnaryInfC{$!B \lor !C$}
\AxiomC{$\Discharge{!B}{x}$}
\RightLabel{$\delta_3$}
\DeduceC{$!D$}
\AxiomC{$\Discharge{!C}{y}$}
\RightLabel{$\delta_4$}
\DeduceC{$!D$}
\DischargeRule{\Elim{\lor}}{x\,y}
\TrinaryInfC{$!D$}
\DisplayProof
\quad\text{بـ}\quad
\bottomAlignProof
\AxiomC{}
\RightLabel{$\delta_2$}
\DeduceC{$\lfalse$}
\RightLabel{\FalseInt}
\UnaryInfC{$!D$}
\DisplayProof
\]
لربما أنشأنا قطعًا جديدًا~!!{formula} قطعه~$!D$، ولا ضمان أن
$\depth{!D} < \depth{!B \lor !C}$!{} لذلك ينبغي أن نسلك هنا المسلك نفسه
الذي سلكناه في اختزال~\Intro{\lor}/\Elim{\lor}، فنستبدل به ما يأتي:
\[
\bottomAlignProof
\AxiomC{}
\RightLabel{$\delta_2$}
\DeduceC{$\lfalse$}
\RightLabel{\FalseInt}
\UnaryInfC{$!B$}
\RightLabel{$\delta_3$}
\DeduceC{$!D$}
\DisplayProof
\quad\text{أو}\quad
\bottomAlignProof
\AxiomC{}
\RightLabel{$\delta_2$}
\DeduceC{$\lfalse$}
\RightLabel{\FalseInt}
\UnaryInfC{$!C$}
\RightLabel{$\delta_4$}
\DeduceC{$!D$}
\DisplayProof
\]

\end{document}
